\chapter*{ABSTRAK}
\addcontentsline{toc}{chapter}{ABSTRAK}

\begin{center}

    \textbf{Implementasi \textit{Graph Retrieval-Augmented Generation} untuk Meningkatkan Presisi \textit{Retrieval} dan Validitas Sintaksis pada Konfigurasi Kubernetes} \\
    \vspace{0.4cm}
    oleh \\
    \vspace{0.4cm}
    Jihan Aurelia\\
    18222001\\
\end{center}


\begin{spacing}{1.0}
Kubernetes menjadi standar orkestrasi kontainer \textit{cloud-native}, namun kompleksitas konfigurasi YAML mendorong penggunaan \textit{Large Language Model} (LLM) yang rentan menghasilkan halusinasi berupa konfigurasi yang keliru secara semantik atau tidak valid secara sintaksis. Pendekatan \textit{Retrieval-Augmented Generation} (RAG) berbasis vektor terbukti tidak memadai karena gagal menangkap relasi hierarkis antarkomponen Kubernetes sehingga kualitas \textit{retrieval} struktural tetap rendah.

\begin{sloppypar}
Penelitian ini membangun sistem \textit{Graph Retrieval-\allowbreak{}Augmented Generation} (GraphRAG) dari spesifikasi OpenAPI Kubernetes v1.30, menghasilkan \textit{knowledge graph} dengan 725 \textit{node}, 18 jenis \textit{edge}, dan 7 kategori relasi. Mekanisme \textit{retrieval} menggabungkan pencocokan nama eksak, \textit{fallback} kemiripan vektor, dan \textit{multi-hop graph traversal} dengan kedalaman adaptif berbasis \textit{intent} ($d = 2$ atau $d = 3$). Validasi YAML diimplementasikan dalam tiga lapis berbasis \textit{knowledge graph}. \textit{Pipeline} dibangun di atas LangGraph menggunakan GPT-4o-mini sebagai \textit{thinker} dan \textit{speaker}, Neo4j sebagai graf dan penyimpanan vektor, serta SQLite sebagai memori percakapan.
\end{sloppypar}

Penelitian ini mengevaluasi GraphRAG terhadap Vector RAG dan Vanilla LLM menggunakan 102 \textit{fixture} tervalidasi pakar. GraphRAG unggul signifikan pada dimensi \textit{retrieval} (RetQ F1 0,7089), jauh melampaui kedua sistem pembanding yang gagal menangkap struktur relasional. Pada dimensi penalaran, \textit{Hop Accuracy} mencapai 0,7562 dengan akurasi lebih tinggi pada kueri berjangkauan sempit dibandingkan kueri lintas-relasi luas, sementara \textit{Faithfulness} sebesar 0,3055 unggul atas Vector RAG meski sebagian klaim jawaban masih bersumber dari pengetahuan parametrik LLM alih-alih graf pengetahuan. Kualitas jawaban akhir (AnsQ) kedua sistem relatif setara secara statistik. Seluruh keunggulan dikonfirmasi signifikan melalui uji Wilcoxon dan \textit{bootstrap} dengan koreksi Holm-Bonferroni.

Temuan ini menegaskan bahwa kontribusi utama GraphRAG terletak pada ketertelusuran relasional, bukan pada kefasihan jawaban semata.

Kata kunci: \textit{Graph Retrieval-Augmented Generation}, \textit{knowledge graph}, Kubernetes, YAML, halusinasi LLM, \textit{intent-adaptive depth traversal}, \textit{faithfulness}, \textit{hop accuracy}.
\end{spacing}
